\subsection{正规子群胚；群胚的商}

定义 7. —— 设 G 为群胚。称 G 的子群胚 H 为正规子群胚，如果 Som(H) = Som(G)，并且对 H 的顶点的每一对 (a, b) 及每条箭 \(f \in \mathrm{G}_{a,b}\) ，有 \(Int(f)(H_b) = H_a\) 。

设 G 为群胚，H 为 G 的子群胚，其顶点集等于 Som(G)。为验证 H 是正规子群胚，只需证明 \(fH_{b}f^{-1} \subset H_{a}\) 对一切 \(a \in \text{Som}(G)\) 、一切 \(b \in \text{Som}(\mathbf{G})\) 以及每条 \(f \in \mathbf{G}_{a,b}\) 皆然。事实上，若此成立，则也有 \(f^{-1}\mathrm{H}_a f \subset \mathrm{H}_b\) （对每条箭 \(f \in \mathbf{G}_{b,a}\) ），即 \(\mathrm{H}_a \subset f\mathrm{H}_b f^{-1}\) ，从而 \(f\mathrm{H}_b f^{-1} = \mathrm{H}_a\) 。

设 G 为群胚，H 为 G 的正规子群胚。对 G 的每个顶点 a，子群 \(H_{a}\) （\(G_{a}\) 的子群）是 \(G_{a}\) 的正规子群。

设 \(\varphi\colon\mathrm{G}\to\mathrm{G}^{\prime}\) 为群胚态射。\(\varphi\) 的核（II, 第 166 页, 例 5）是 G 的正规子群胚。事实上，设 \(f\) 为 G 中连接 \(a\) 与 \(b\) 的箭，并设 \(g\in\mathrm{Ker}(\varphi)_{b}\) ；则我们有 \(\varphi(fgf^{-1})=\varphi(f)\varphi(g)\varphi(f^{-1})=\varphi(f)e_{\varphi(b)}\varphi(f)^{-1}=e_{\varphi(a)}\) ，由此得到包含关系 \(f\mathrm{Ker}(\varphi)_{b}f^{-1}\subset\mathrm{Ker}(\varphi)_{a}\) 。

设 G 为群胚。群胚 G 以及箭集为 G 的单位元全体的那个 G 的子群胚，都是 G 的正规子群胚。G 的一族正规子群胚的交是 G 的正规子群胚。特别地，对每个子集 F ⊂ Fl(G)，存在箭集包含 F 的最小的 G 的正规子群胚。称之为由 F 生成的正规子群胚。

设 H 为 G 的正规子群胚。设 R 为 \(\mathrm{Fl}(\mathrm{G})\) 中由“ \(R\{f,g\}\) 当且仅当存在 \(\mathrm{Fl}(\mathrm{H})\) 中的箭 x 与 y 使得 f = xgy”定义的等价关系。若 f 与 g 是模 R 等价的 G 的箭，则它们的起点（相应地，终点）属于 H 的同一轨道。令 \(F' = \mathrm{Fl}(\mathrm{G}) / \mathcal{R}\) ，以 \(o'\) 与 \(t'\) 表示由 o 与 t 通过取商而得的从 \(F'\) 到 \(\mathrm{Orb}(\mathrm{H})\) 的映射。以 G/H 表示箭图 \((\mathrm{Orb}(\mathrm{H}), \mathrm{F}', o', t')\) 。

引理. —— 对 G/H 的给定的可合成箭 u 与 v，可以在 Fl(G) 中选取 u 与 v 的可合成的代表 f 与 g。积 fg 模 R 的类不依赖于 f 与 g 的选取。

设 \(f\) 与 \(g\) 为 \(u\) 与 \(v\) 在 \(\mathrm{Fl(G)}\) 中的任意代表。\(f\) 的终点与 \(g\) 的起点属于 H 的同一轨道，因而可用 H 的一条箭 \(x\) 连接（II, 第 162 页, 命题 1）。于是箭 \(fx\) 与 \(g\) 是 \(u\) 与 \(v\) 在 F 中的在 G 内可合成的代表。这就证明了第一个断言。

现在设 \(f\) 、\(g\) 与 \(f'\) 、\(g'\) 分别为 \(u\) 与 \(v\) 的在 G 内可合成的代表。按假设，存在 H 中的箭 \(x\) 、\(y\) 、\(z\) 、\(t\) ，使得 \(f' = xfy\) 且 \(g' = zgt\) 。于是有 \(yz \in \mathrm{H}_{t(f)}\) 及 \(f'g' = x f y z g t = x f(yz)f^{-1}f g t\) 。由于 H 是 G 的正规子群胚，箭 \(f(yz)f^{-1}\) 是 H 的箭。因此箭 \(x f(yz)f^{-1}\) 亦然，这就证明了 \(f'g'\) 与 \(fg\) 模 \(\mathcal{R}\) 等价。

根据此引理，在箭图 G/H 中存在唯一的合成法则 \(m\) ，使得对每一对可合成箭 \((u,v)\) ，\(m(u,v)\) 都是模 \(\mathcal{R}\) 由积 \(fg\) 所代表的类，其中 \((f,g)\) 取遍 G 中满足 \(u\) 为 \(f\) 的类且 \(v\) 为 \(g\) 的类的可合成箭对。赋有此合成法则的 G/H 是一个群胚。设 \(p_1\colon \operatorname{Som}(\mathrm{G}) \to \operatorname{Som}(\mathrm{G}/\mathrm{H})\) 与 \(p_2\colon \operatorname{Fl}(\mathrm{G}) \to \operatorname{Fl}(\mathrm{G}/\mathrm{H})\) 为典范满射。偶 \(p = (p_1,p_2)\) 是从 G 到 G/H 的群胚态射，其核为正规子群胚 H。

定义 8. —— 我们称 G/H 为 G 对 H 的商群胚，称 p: G → G/H 为典范态射。

注. —— 1) 映射 Som(p) 通过取商定义从 G 的轨道集到 G/H 的轨道集上的一个双射。特别地，G 传递当且仅当 G/H 传递。

\begin{enumerate}
\def\labelenumi{\arabic{enumi})}
\setcounter{enumi}{1}
\tightlist
\item
  设 \(a\) 与 \(b\) 为 G 的顶点。从 \(\mathrm{G}_{a,b}\) 到 \((\mathrm{G} / \mathrm{H})_{p(a),p(b)}\) 的、由 \(p\) 诱导的映射是满射。事实上，设 \(u\) 为 G/H 中连接 \(p(a)\) 与 \(p(b)\) 的一条箭；它是 G 的某条箭的模 \(\mathcal{R}\) 类，设该箭为 \(f\) 。其起点 \(o(f)\) （\(f\) 的起点）属于 \(a\) 在 H 中的轨道；因此存在 H 中的一条箭 \(x\) 连接 \(a\) 与 \(o(f)\) 。类似地，存在 H 中的一条箭 \(y\) 连接 \(t(f)\) 与 \(b\) 。于是，\(f' = xfy\) 是 \(\mathrm{G}_{a,b}\) 的元素，其模 \(\mathcal{R}\) 的类为 \(u\) 。
\end{enumerate}

命题 2. —— 设 a 为 G 的顶点。由 p 通过取迷向群而得的群同态 \(p_{a}\colon\mathrm{G}_{a}\to(\mathrm{G}/\mathrm{H})_{p(a)}\) 是满射的；其核为 \(\mathrm{H}_{a}\) 。

由前面的注，同态 \(p_a\) 是满射的。其核包含 \(\mathrm{H}_a\) ，这由群胚 \(\mathrm{G / H}\) 的构造可知。设 \(f\) 为 \(\mathrm{G}_a\) 中使得 \(p_a(f) = e_{p(a)}\) 的元素。这意味着 \(f\) 与 \(e_a\) 模 \(\mathcal{R}\) 等价；于是存在 H 的箭 \(x\) 与 \(y\) ，使得 \(f = xe_a y\) 。必然地，\(x\) 与 \(y\) 属于 \(\mathrm{H}_a\) ，由此 \(f \in \mathrm{H}_a\) 。

命题 3. —— 设 G 为群胚，H 为 G 的正规子群胚，p: G → G/H 为典范态射。设 \(\varphi\colon G\to G'\) 为群胚态射，满足 \(H\subset\operatorname{Ker}(\varphi)\) 。则存在唯一的群胚态射 \(\overline{\varphi}\colon G/H\to G'\) ，使得 \(\overline{\varphi}\circ p=\varphi\) 。

我们称 \(\overline{\varphi}\) 为 \(\varphi\) 通过取商诱导的群胚态射。

这种态射的唯一性是显然的，因为映射 \(\operatorname{Som}(p)\) 与 \(\operatorname{Fl}(p)\) 都是满射。

设 \(a\) 与 \(b\) 为 G 的顶点。若它们属于 H 的同一轨道，则存在 H 中的一条箭 \(f\) 连接 \(a\) 与 \(b\) ，且有 \(\varphi(f) = e_{\varphi(a)}\) 。特别地，\(\varphi(a) = \varphi(b)\) 。于是，映射 \(\operatorname{Som}(\varphi)\) 通过取商定义了一个映射 \(\overline{\varphi}_1\colon \operatorname{Orb}(\mathrm{H}) \to \operatorname{Som}(\mathrm{G}')\) 。设 \(f\) 与 \(g\) 为 G 中的箭。若 \(f\) 与 \(g\) 模 \(\mathcal{R}\) 等价，则存在 H 中的箭 \(x\) 与 \(y\) ，使得 \(f = xgy\) 。于是，\(\varphi(f) = \varphi(x)\varphi(g)\varphi(y) = \varphi(g)\) ，因为 \(\varphi(x)\) 与 \(\varphi(y)\) 是单位元。因此，映射 \(\operatorname{Fl}(\varphi)\) 通过取商定义了一个映射 \(\overline{\varphi}_2\colon \operatorname{Fl}(\mathrm{G}) / \mathcal{R} \to \operatorname{Fl}(\mathrm{G}')\) 。

设 \(f\) 与 \(g\) 为 \(G\) 的可合成箭；以 \(u\) 与 \(v\) 表示它们在 \(\mathrm{Fl}(\mathrm{G} / \mathrm{H})\) 中的类。我们有 \(\overline{\varphi}_2(uv) = \varphi(fg) = \varphi(f)\varphi(g) = \overline{\varphi}_2(u)\overline{\varphi}_2(v)\) 。

偶 \(\overline{\varphi} = (\overline{\varphi}_1, \overline{\varphi}_2)\) 是从 G/H 到 G' 的群胚态射，且有 \(\overline{\varphi} \circ p = \varphi\) 。

